\paragraph{Layout-independent norms.} To test whether the LLM proposes domain-level norms rather than layout-specific
rules, we built two additional stores for the spatial domain: a second training layout with a horizontal aisle, and
a test layout (never shown to the LLM) whose aisles lie where the training layouts had ordinary free cells. The
candidate set always contained coordinate-based competitors that explain the training demonstrations equally well.
Without guidance, none of 10 GPT-5.5 refinements contained a norm equivalent to the intended one across all
layouts. Its relational proposals (e.g.\ ``between shelves on the east and west'') matched the intended norm only on
the training layout. Adding a short instruction to separate domain-level norms from rules about the particular
layout, without naming the concept, removed all coordinate-based proposals and produced a layout-equivalent norm in
6/10 runs. However, with a single training layout the Bayesian learner still selected the shorter coordinate rule in
every run, and test-layout accuracy stayed at chance (0.50, violation F1 0.00). With demonstrations from two
layouts, coordinate rules were never selected: the intended norm was selected in 5/10 runs and test-layout accuracy
rose to 0.90 (F1 0.92). The remaining errors came from relational but over-general definitions of the concept.
Proposal, selection and generalisation are therefore distinct: guidance improves what the LLM proposes, but
diverse environments are needed for the evidence to favour the generalisable norm.
